\chapter{CA 19-9 (Cancer Antigen 19-9)}

\section*{What Is This Test?}
CA 19-9 (carbohydrate antigen 19-9) is a tumor-associated carbohydrate antigen related to the sialylated Lewis-a (sLe$^\text{a}$) blood group antigen. It is a modified form of the Lewis-a blood group antigen that is overexpressed on the surface of cancer cells, particularly pancreatic adenocarcinoma, biliary tract cancer (cholangiocarcinoma), gastric cancer, and colorectal cancer. CA 19-9 is synthesized by exocrine epithelial cells of the pancreas, biliary tract, gastric mucosa, and salivary glands. In pancreatic ductal adenocarcinoma, CA 19-9 is the most clinically useful serum tumor marker for monitoring treatment response and detecting recurrence. However, its use as a screening test is limited by inadequate sensitivity (elevated in only 70--80\% of pancreatic cancers) and poor specificity (elevated in benign hepatobiliary and pancreatic conditions including cholangitis, pancreatitis, and obstructive jaundice). Importantly, approximately 5--10\% of the population are Lewis-a negative (Le$^\text{a-b-}$) and are genetically incapable of synthesizing CA 19-9, so they will have undetectable CA 19-9 even with advanced pancreatic cancer.

\section*{Alternative Names}
CA 19-9; Cancer Antigen 19-9; Carbohydrate Antigen 19-9; Sialylated Lewis-a Antigen; sLe$^\text{a}$; CA 19-9 Radioimmunoassay; GI Cancer Antigen; Pancreatic Cancer Marker

\section*{Principle}
CA 19-9 is measured by a two-site sandwich chemiluminescent or electrochemiluminescent immunoassay using the monoclonal antibody 1116 NS 19-9, which was generated by immunizing mice with the SW1116 human colorectal carcinoma cell line. The antibody recognizes the sialylated Lewis-a (sLe$^\text{a}$) epitope. In the assay, patient serum is incubated with a capture antibody and a labeled detection antibody (both recognizing CA 19-9). After incubation, the immune complex is captured, washed, and signal is generated by chemiluminescence (ruthenium-label ECLIA on Roche Cobas, acridinium-ester CLIA on Abbott Architect, or similar platforms). Signal intensity is proportional to CA 19-9 concentration. Results are reported in U/mL. CA 19-9 is calibrated against an internal reference standard; there is no international standard. Inter-method variation between manufacturers is significant; serial CA 19-9 measurements must be performed using the same assay. Heterophile antibodies and human anti-mouse antibodies (HAMA) can interfere with immunoassay detection and cause falsely elevated results. High-dose hook effect (antigen excess producing falsely low results) should be considered in patients with massive tumor burdens: the laboratory may perform dilution studies if CA 19-9 is unexpectedly normal in a patient with known metastatic disease.

\section*{History}
CA 19-9 was discovered in 1979 by Hilary Koprowski and colleagues at the Wistar Institute. After immunizing mice with the SW1116 colorectal cancer cell line, they identified monoclonal antibody 1116 NS 19-9, which reacted with a novel carbohydrate antigen. The target epitope was identified as sialylated Lewis-a (sLe$^\text{a}$), a modified blood group antigen whose synthesis requires the Lewis enzyme fucosyltransferase (encoded by the FUT3 gene). The CA 19-9 assay was introduced into clinical practice in the early 1980s. Magnani et al. (1982) demonstrated that CA 19-9 is the ligand for E-selectin, implicating it in the hematogenous metastasis of cancer cells. CA 19-9 is now the most widely used serum tumor marker for pancreatic cancer, with its primary role in monitoring treatment response and detecting recurrence rather than screening. The National Comprehensive Cancer Network (NCCN) and American Society of Clinical Oncology (ASCO) recommend against CA 19-9 for screening but endorse its use for monitoring.

\section*{Why This Test Is Ordered}
\begin{itemize}
  \item Monitoring treatment response in patients with known pancreatic ductal adenocarcinoma --- falling CA 19-9 levels after surgical resection or during chemotherapy indicate response; rising levels suggest progression or recurrence
  \item Preoperative CA 19-9 measurement for prognostication in potentially resectable pancreatic cancer --- a CA 19-9 level >500 U/mL is associated with a lower likelihood of R0 (complete) resection and poorer overall survival, and some centers use this as a relative contraindication to primary surgery (preferring neoadjuvant chemotherapy)
  \item Detection of recurrence after pancreatic cancer surgery --- rising CA 19-9 typically precedes radiographic evidence of recurrence by 2--6 months
  \item Diagnosis and monitoring of cholangiocarcinoma (bile duct cancer), where CA 19-9 is elevated in 60--80\% of patients
  \item Monitoring of gastric cancer and colorectal cancer in patients with known CA 19-9 elevation that correlates with tumor burden
  \item Differentiation of benign biliary strictures (primary sclerosing cholangitis, chronic pancreatitis) from malignant strictures (cholangiocarcinoma) when combined with imaging and cytology. CA 19-9 >100 U/mL raises concern for malignancy, though marked elevations (>1,000 U/mL) can occur in benign obstructive jaundice without cancer
  \item Adjunct to imaging in the evaluation of suspected pancreatic cancer in patients presenting with painless jaundice, epigastric pain radiating to the back, unexplained weight loss, and new-onset diabetes
\end{itemize}

\section*{Is This a Primary or Secondary Test}
CA 19-9 is a secondary test. It is NOT recommended as a screening test for pancreatic cancer in asymptomatic individuals due to low sensitivity (70--80\%) and poor specificity (elevated in acute and chronic pancreatitis, cholangitis, obstructive jaundice). CA 19-9 is used primarily for monitoring treatment response and detecting recurrence in patients with confirmed pancreaticobiliary malignancies.

\section*{How This Test Is Performed}
For most patients, a health care professional will take a blood sample from a vein in your arm using a small needle. After the needle is inserted, a small amount of blood will be collected into a test tube or vial. You may feel a little sting when the needle goes in or out. This usually takes less than five minutes.

\section*{What Preparation Is Needed}
No special preparation is required. Fasting is not necessary. CA 19-9 should ideally be measured when the patient is not acutely ill with cholangitis, acute pancreatitis, or biliary obstruction, as these conditions can cause marked transient elevations. For baseline CA 19-9 before pancreatic surgery, the measurement should be obtained after relief of biliary obstruction (by endoscopic stent or percutaneous drain) and normalization of bilirubin, as obstructive jaundice can artificially elevate CA 19-9 by an order of magnitude.

\section*{Contraindications and Precautions}
No absolute contraindications. The most critical limitation of CA 19-9 testing is that approximately 5--10\% of individuals of European descent (and up to 20\% of individuals of African descent) have the Lewis-a-negative (Le$^\text{a-b-}$) blood group phenotype and are genetically unable to synthesize sialylated Lewis-a antigen, the epitope recognized by the CA 19-9 assay. In these individuals, CA 19-9 will be undetectable even with advanced pancreatic cancer. Lewis blood group testing is rarely performed in clinical practice, but a CA 19-9 of zero or near-zero in a patient with known pancreatic cancer suggests Lewis-null status and renders CA 19-9 useless for monitoring. CA 19-9 is elevated in a wide variety of benign hepatobiliary and pancreatic conditions: acute and chronic pancreatitis, acute cholangitis, obstructive jaundice (any cause --- gallstones, benign stricture), cirrhosis, hepatitis, and cystic fibrosis. CA 19-9 levels correlate closely with total bilirubin in biliary obstruction: levels >1,000 U/mL are common during cholangitis and acute biliary obstruction, and these levels decrease dramatically (often into the normal range) after relief of obstruction by stenting. Other malignancies that elevate CA 19-9: gastric adenocarcinoma, colorectal adenocarcinoma, hepatocellular carcinoma, cholangiocarcinoma, gallbladder carcinoma, and occasionally ovarian cancer (particularly mucinous type) and lung cancer. Specimen should be collected in a serum separator tube (SST, gold-top) or red-top tube. Heterophile antibodies and HAMA can produce false-positive elevations. High-dose hook effect (antigen excess) can produce falsely low results in massive tumor burden; the laboratory may need to dilute the sample.

\section*{Risks}
Minimal risk. Slight pain or bruising at the venipuncture site. Rare: hematoma, infection, vasovagal syncope.

\section*{What Happens After the Test}
Results are typically available within 1 to 2 hours. An elevated CA 19-9 in a patient with a pancreatic mass or biliary stricture supports the diagnosis of malignancy but is not diagnostic in itself. CA 19-9 should be interpreted alongside cross-sectional imaging (multiphasic CT pancreas protocol or MRI/MRCP), endoscopic ultrasound (EUS) with fine-needle aspiration (FNA) for tissue confirmation, and the patient's clinical presentation. A CA 19-9 >1,000 U/mL in a jaundiced patient should be repeated 2--4 weeks after biliary decompression to determine if the elevation is due to obstructive jaundice alone or reflects underlying malignancy.

\section*{Understanding Results}

\subsection*{Below Normal}
\begin{itemize}
  \item CA 19-9 <37 U/mL: Normal in healthy individuals. Does not exclude pancreatic cancer --- 20--30\% of patients with pancreatic adenocarcinoma (exocrine) have normal CA 19-9 at diagnosis, particularly small tumors (<2 cm) and patients with early-stage resectable disease
  \item Lewis-a-negative (Le$^\text{a-b-}$) phenotype: These individuals (5--10\% of population) cannot synthesize the sialylated Lewis-a antigen due to a mutation in the FUT3 fucosyltransferase gene. CA 19-9 is undetectable even with advanced pancreatic adenocarcinoma. If a patient with known pancreatic cancer has a consistently undetectable CA 19-9, CA 19-9 cannot be used for monitoring. CEA and other markers (CA 125) may serve as alternative monitoring tools in these patients
  \item Normalization of CA 19-9 after curative-intent surgery for pancreatic cancer: Indicates complete resection (R0) and favorable prognosis. Persistently elevated CA 19-9 after surgery suggests residual microscopic disease and is associated with early recurrence
  \item Normalization of CA 19-9 during chemotherapy for advanced pancreatic cancer: Indicates treatment response. A decrease of $\geq$50\% in CA 19-9 during chemotherapy is associated with improved progression-free and overall survival
  \item Isolated normal CA 19-9 with elevated bilirubin without malignancy: The absence of CA 19-9 elevation in a jaundiced patient does not exclude pancreatic cancer but argues against a large primary pancreatic tumor
\end{itemize}

\subsection*{Above Normal}
\begin{itemize}
  \item CA 19-9 37--100 U/mL: Mild elevation. Non-specific --- can be due to benign pancreatic or biliary conditions (chronic pancreatitis, mild cholangitis, benign biliary stricture), non-pancreatic malignancies (gastric, colorectal, biliary), or early-stage pancreatic cancer. Warrants clinical correlation with imaging
  \item CA 19-9 100--500 U/mL: Moderate elevation. Increases suspicion for pancreatic adenocarcinoma. In the appropriate clinical setting (painless jaundice, weight loss, new-onset diabetes), warrants urgent cross-sectional imaging (CT pancreas protocol or MRI/MRCP) and EUS with FNA
  \item CA 19-9 >500 U/mL: High probability of pancreaticobiliary malignancy. The preoperative CA 19-9 level predicts resectability --- CA 19-9 >500 U/mL is associated with a higher likelihood of locally advanced or metastatic disease on surgical exploration and incomplete (R1/R2) resection
  \item CA 19-9 >1,000 U/mL: Strongly suspicious for pancreatic adenocarcinoma. In the presence of obstructive jaundice, this level may reflect biliary obstruction amplifying CA 19-9 rather than massive tumor; repeat after biliary decompression. In the absence of jaundice, >1,000 U/mL is strongly suggestive of advanced or metastatic disease
  \item CA 19-9 >10,000--100,000 U/mL: Seen in advanced metastatic pancreatic cancer with large tumor burden
  \item Rising CA 19-9 after surgery or during chemotherapy: Indicates disease progression or recurrence. Rising CA 19-9 typically precedes radiographic evidence of recurrence by 2--6 months. The doubling time of CA 19-9 is a negative prognostic factor
  \item CA 19-9 elevation from benign conditions: Acute cholangitis (CA 19-9 can reach >10,000 U/mL during active bacterial cholangitis and falls rapidly with antibiotic therapy and biliary drainage), acute pancreatitis (moderate to marked elevation during the acute episode, returning to normal over 4--6 weeks), chronic pancreatitis (mild to moderate elevation, <500 U/mL), obstructive jaundice from benign causes (choledocholithiasis, benign biliary stricture, primary sclerosing cholangitis --- CA 19-9 is often elevated in PSC even without superimposed cholangiocarcinoma)
  \item Non-pancreatic malignancies: Cholangiocarcinoma (the second most common indication for CA 19-9), gastric adenocarcinoma, colorectal adenocarcinoma, hepatocellular carcinoma, gallbladder carcinoma, and ovarian mucinous carcinoma
\end{itemize}

\section*{Normal Results}
\begin{itemize}
  \item \textbf{CA 19-9:} <37 U/mL (upper limit of normal for most laboratories; some use <35 U/mL)
  \item \textbf{Borderline:} 37--55 U/mL (equivocal; clinical significance uncertain)
  \item \textbf{Mild elevation:} 55--100 U/mL
  \item \textbf{Moderate elevation:} 100--500 U/mL
  \item \textbf{Marked elevation:} >500 U/mL
  \item \textbf{Lewis-a-negative individuals (Le$^\text{a-b-}$):} CA 19-9 is undetectable or near-zero (<2 U/mL) even with advanced pancreatic cancer. Affects approximately 5--10\% of the population
  \item CA 19-9 results in a jaundiced patient should be interpreted cautiously, and the test should be repeated after relief of biliary obstruction. A CA 19-9 that drops by >90\% after stenting strongly suggests that the elevation was due to obstruction rather than tumor
  \item Serial CA 19-9 measurements should be performed using the same assay method to ensure comparability
\end{itemize}

\section*{Follow-up Tests}
\begin{itemize}
  \item \textbf{Multiphasic CT of the abdomen (pancreas protocol):} The first-line imaging modality for suspected pancreatic cancer. Includes non-contrast, arterial phase (for arterial anatomy), pancreatic parenchymal phase (for pancreatic mass enhancement), and portal venous phase (for liver metastases and venous anatomy). The mass typically appears as a hypodense (hypovascular) lesion relative to the normally enhancing pancreatic parenchyma
  \item \textbf{MRI/MRCP (magnetic resonance cholangiopancreatography):} More sensitive than CT for small pancreatic tumors and intraductal papillary mucinous neoplasms (IPMN). MRCP provides non-invasive imaging of the pancreatic and biliary ductal system without contrast injection
  \item \textbf{Endoscopic ultrasound (EUS) with fine-needle aspiration (FNA) or fine-needle biopsy (FNB):} The most sensitive modality for detecting small pancreatic tumors (<2 cm). Provides tissue for cytology and histologic confirmation of malignancy. EUS-guided celiac plexus neurolysis can be performed for pain palliation
  \item \textbf{ERCP (endoscopic retrograde cholangiopancreatography) with biliary brushing and stenting:} For patients with obstructive jaundice. Provides biliary decompression (relieves pruritus, reduces the risk of cholangitis, permits normalization of liver function before surgery or chemotherapy) and allows biliary cytology sampling
  \item \textbf{CEA (carcinoembryonic antigen):} An alternative tumor marker for pancreatic cancer monitoring, particularly in CA 19-9 non-secretors (Lewis-null phenotype). CEA is elevated in 40--60\% of pancreatic cancers
  \item \textbf{CA 125:} Elevated in a subset of pancreatic cancers and may serve as an alternative monitoring marker in Lewis-null individuals or patients with mucinous pancreatic tumors
  \item \textbf{Liver function tests (AST, ALT, ALP, total and direct bilirubin):} To assess the degree of biliary obstruction and hepatic function. Obstructive jaundice shows a predominantly direct (conjugated) hyperbilirubinemia with markedly elevated alkaline phosphatase and GGT
  \item \textbf{CA 19-9 after biliary decompression:} When CA 19-9 is elevated in the setting of obstructive jaundice, the test should be repeated 3--4 weeks after successful biliary stenting or drainage, when bilirubin has normalized. The post-decompression CA 19-9 level is the appropriate value for prognostication and treatment monitoring
  \item \textbf{Laparoscopy with peritoneal washings:} In patients with borderline resectable or locally advanced pancreatic cancer, diagnostic laparoscopy can detect occult peritoneal metastases not visualized on CT, avoiding unnecessary laparotomy
\end{itemize}

\section*{References}
\begin{enumerate}
  \item Koprowski H, Steplewski Z, Mitchell K, et al. Colorectal carcinoma antigens detected by hybridoma antibodies. Somatic Cell Genet. 1979;5(6):957--971.
  \item Magnani JL, Brockhaus M, Smith DF, et al. A monosialoganglioside is a monoclonal antibody-defined antigen of colon carcinoma. Science. 1981;212(4490):55--56.
  \item Tempero MA, Malafa MP, Al-Hawary M, et al. Pancreatic Adenocarcinoma, Version 2.2021, NCCN Clinical Practice Guidelines in Oncology. J Natl Compr Canc Netw. 2021;19(4):439--457.
  \item Ballardini M, Reni M, Bazzoli E, et al. CA 19-9 in patients with pancreatic cancer undergoing treatment: a systematic review. Ann Oncol. 2011;22(Suppl 4):iv175.
  \item Ballehaninna UK, Chamberlain RS. Serum CA 19-9 as a biomarker for pancreatic cancer --- a comprehensive review. Indian J Surg Oncol. 2012;3(2):121--127.
\end{enumerate}

\clearpage
